\begin{solution}{94}
\begin{subsolution}
平行激光束照射弹簧可先、后照到两小段倾斜的细丝，沿光束方向看，两段细丝呈现交叉状。根据题给说明，两者分别产生与单缝衍射相同的条纹。通过图 a 右图测量可得，各组单缝衍射条纹中点连线之间的夹角为
\begin{equation}
\phi = 3 0 ^ {\circ}
\label{eq:1.s94}
\end{equation}
细丝直径和两根相邻的平行细丝间距可通过衍射条纹的细节结构计算得到。可以看到，图 a 右图中的衍射条纹有一种长的空间周期和一种短的空间周期，分别对应细丝的衍射条纹和相邻细丝间空隙的衍射条纹。设第 $n$ 级暗纹在屏上的位置为 $x_{n}$ ，由于光屏远离细丝，衍射角 $\theta_{n}$ 可视为很小的角，即对第 $n$ 级暗纹近似有
\begin{equation}
d \sin \theta_ {n} \approx d \theta_ {n} = d \frac {x _ {n}}{L} = n \lambda
\label{eq:2.s94}
\end{equation}
式中 $\lambda$ 是激光的波长。相邻暗纹间距为
\begin{equation}
X \equiv x _ {n + 1} - x _ {n} = \lambda \frac {L}{d}
\label{eq:3.s94}
\end{equation}
通过图 a 右图进行测量，可得到长的空间周期为
\begin{equation}
X _ {1} = 7. 9 \mathrm{mm}
\label{eq:4.s94}
\end{equation}
由\eqref{eq:3.s94}式可得细丝直径为
\begin{equation}
d _ {1} = \lambda \frac {L}{X _ {1}} = 0. 3 4 \mathrm{mm}
\label{eq:5.s94}
\end{equation}
通过图 a 右图进行测量，可得到短的空间周期为
\begin{equation}
X _ {2} = 1. 7 \mathrm{mm}
\label{eq:6.s94}
\end{equation}
由\eqref{eq:3.s94}式可得相邻平行细丝的间距为
\begin{equation}
d _ {2} = \lambda \frac {L}{X _ {2}} = 1. 6 \mathrm{mm}
\label{eq:7.s94}
\end{equation}
如右图所示，螺旋线在平面上展开即成为一直角三角形的斜边，两条直角边的长度分别为螺距 $p$ 和圆周长 $2\pi R$ ， $E$ 点为斜边中点。由相邻平行细丝间距 $d_{2}$ 和两交叉细丝夹角 $\phi$ 可算出螺旋结构的螺距
\begin{equation}
p = \frac {d _ {2}}{\cos \frac {\phi}{2}} = \frac {\lambda L}{X _ {2} \cos \frac {\phi}{2}}
\label{eq:8.s94}
\end{equation}
将已得到的数据（见\eqref{eq:1.s94}\eqref{eq:7.s94}式）代入\eqref{eq:8.s94}式得
\begin{equation}
p = 1. 6 \mathrm{mm}
\label{eq:9.s94}
\end{equation}
由右图中几何关系得
\begin{equation}
\frac {2 \pi R}{p} = \tan \theta = \cot \frac {\phi}{2}
\label{eq:10.s94}
\end{equation}
将\eqref{eq:8.s94}式代入上式得，螺旋结构俯视圆的半径为
\begin{equation}
R = \frac {d _ {2}}{2 \pi \sin \frac {\phi}{2}} = \frac {\lambda L}{2 \pi X _ {2} \sin \frac {\phi}{2}}
\label{eq:11.s94}
\end{equation}
将已得到的数据代入\eqref{eq:11.s94}式得
\begin{equation}
R = 0. 9 6 \mathrm{mm}
\label{eq:12.s94}
\end{equation}
\end{subsolution}
\begin{subsolution}
由图 b 可测出，各组衍射条纹中心的连线之间的夹角 $\phi'$ 为
\begin{equation}
\phi^ {\prime} = 8 4 ^ {\circ}
\label{eq:13.s94}
\end{equation}
测出条纹的空间周期 $X'$ 为
\begin{equation}
X ^ {\prime} = 4. 7 \mathrm{mm}
\label{eq:14.s94}
\end{equation}
利用\eqref{eq:3.s94}式，相邻分子链间距为
\begin{equation}
d ^ {\prime} = \lambda^ {\prime} \frac {L ^ {\prime}}{X ^ {\prime}} = 2. 9 \mathrm{nm}
\label{eq:15.s94}
\end{equation}
根据分子链间距 $d'$ （见\eqref{eq:15.s94}式）和夹角 $\phi'$ （见\eqref{eq:13.s94}式），由类似\eqref{eq:8.s94}式，可算出螺距 $p'$
\begin{equation}
p ^ {\prime} = \frac {d ^ {\prime}}{\cos \frac {\phi^ {\prime}}{2}} = \frac {\lambda^ {\prime} L ^ {\prime}}{X ^ {\prime} \cos \frac {\phi^ {\prime}}{2}} = 3. 9 \mathrm{nm}
\label{eq:16.s94}
\end{equation}
由类似\eqref{eq:11.s94}式，可算出俯视半径 $R'$
\begin{equation}
R ^ {\prime} = \frac {d ^ {\prime}}{2 \pi \sin \frac {\phi^ {\prime}}{2}} = \frac {\lambda^ {\prime} L ^ {\prime}}{2 \pi X ^ {\prime} \sin \frac {\phi^ {\prime}}{2}} = 0. 6 8 \mathrm{nm}
\label{eq:17.s94}
\end{equation}
\end{subsolution}
\end{solution}